%% Chapter 1: Introduction (chapters/chapter01.tex)

% ==============================================================================
% CONTENT BLUEPRINT — Chapter 1: Introduction
% ==============================================================================
% Reading audience: committee, external examiner, fellow students.
% Goal: deliver a complete 3-pager that frames the project, defends its
% existence, and tells the reader exactly what the rest of the report does.
%
% Recommended length: 8–12 pages. Sections below are required unless marked
% (optional). Use \ref{ch:methodology}, \ref{ch:results},
% \ref{ch:discussion}, \ref{ch:conclusion} (defined in chapters 2–5).
%
% ----------------------------------------------------------------------------
% §1.1 Background and Motivation  (1–2 pages)
%   • Open with the broad field context (e.g. non-linear optics, semiconductor
%     physics, materials science). Define the sub-field in one paragraph.
%   • State why the field matters now — applications, recent breakthroughs,
%     technoeconomic relevance.
%   • Narrow: which sub-problem the project touches (e.g. third-order
%     non-linear refraction, bandgap tuning under strain).
%   • End with a strong "however, X is unresolved / poorly quantified / under-
%     studied at this level" sentence. That sets up §1.2 and §1.3.
%   • Cite 4–6 references in this section — review papers, foundational
%     textbooks, and one or two recent (≤5 yr) articles.
%
% §1.2 Literature Review  (2–3 pages)
%   • Organize chronologically OR thematically — pick ONE and be explicit.
%   • Group studies into 2–3 clusters (foundational theory → measurement
%     techniques → recent materials/applications).
%   • For each cluster: what was done, how, key finding, citation.
%   • End each cluster with a one-line critique ("limited to ...", "did not
%     address ...", "scope restricted to ...").
%   • Final paragraph: explicitly state the gap. This paragraph is the bridge
%     into §1.3.
%   • Cite 8–15 references total. Use \cite{key1,key2} for grouped, distinct
%     entries for distinct claims.
%
% §1.3 Statement of the Problem  (1 paragraph)
%   • One crisp sentence stating what the project is going to resolve, in the
%     format: "Despite [evidence of progress in §1.2], [specific question]
%     remains [unquantified / unresolved / unverified] under [specific
%     conditions]."
%   • Optional: a second sentence quantifying the gap ("...with no published
%     value to within ±20 % at room temperature for X-class materials").
%
% §1.4 Project Objectives  (1 page, bulleted)
%   • 3–6 measurable objectives.
%   • Each bullet must be (a) specific, (b) testable, (c) deliverable by end
%     of project. Bad example: "To study non-linear optics." Good example:
%     "To measure the closed-aperture Z-scan signature of SF59 glass at
%     λ = 632.8 nm and extract n₂ to within ±15 %."
%
% §1.5 Scope and Limitations  (½ page)
%   • What is included (materials, parameters, instruments, time window).
%   • What is excluded (temperature dependence, other non-linear orders,
%     different wavelength ranges, etc.). State exclusions explicitly; do
%     not let the reader guess.
%   • Optional: brief paragraph on assumptions underpinning the work
%     (isotropic media, CW vs pulsed laser, ideal Gaussian beam).
%
% §1.6 Organization of the Report  (½ page)
%   • One short paragraph per subsequent chapter. Tell the reader exactly
%     what to expect so they can decide whether to read sequentially or skip.
%   • End with a forward pointer to §1 or §2 of Chapter 2 where the theory
%     begins.
%
% Writing-style rules for this chapter:
%   • Past tense for what has been reported by others; present tense for
%     what you describe as current understanding.
%   • No first-person plural ("we believe") — use "this project investigates"
%     or "the present work". Exception: §1.4 and §1.5 may use "we measure" /
%     "we compute" since you are describing your own goals.
%   • Keep paragraphs 5–8 lines. Anything longer needs a subsection.
%   • Every equation must be numbered and referenced in discussion (where it
%     is used), not just defined.
% ==============================================================================

\chapter{Introduction}
\label{ch:introduction}

Non-linear optics (NLO) forms the foundation of modern high-speed communication systems, laser technologies, and advanced spectroscopy. While linear optical materials exhibit response fields that scale proportionally with input light field amplitudes, non-linear materials showcase phenomena where the refractive index and absorption profiles change as a function of optical intensity. This chapter introduces the core motivations for investigating light-matter interactions under intense coherent radiation.

\section{Background and Motivation}
In recent years, the need for all-optical switching devices and high-speed modulators has driven intensive research into materials exhibiting high third-order non-linear optical properties. Understanding how intense laser fields interact with novel semiconductor structures is crucial for next-generation telecommunication applications.

\section{Literature Review}
Early investigations of third-order optical effects were pioneered by Boyd \cite{boyd2020nonlinear}. The development of the Z-scan technique by Sheik-Bahae et al. \cite{sheik1990sensitive} provided a simple, sensitive method to resolve both the sign and magnitude of $n_2$. Since then, numerous studies have explored various semiconductor materials for switching applications \cite{sutherland2003handbook}. Recent theoretical extensions have attempted to map these effects using multi-scale simulations \cite{yamada2018numerical}.

Despite these advances, a persistent gap remains in matching experimental data with theoretical simulation schemes for student-level projects, where formatting inconsistencies and missing standardization often limit reproducibility.

\section{Statement of the Problem}
Although many materials exhibit strong NLO properties, their performance limits under continuous-wave vs. pulsed laser exposure are not fully quantified. Furthermore, matching experimental data with theoretical simulation schemes remains a challenge for student projects due to formatting inconsistencies and missing standardization.

\section{Project Objectives}
The core objectives of this B.Sc. project are:
\begin{itemize}
    \item To set up a reliable Z-scan system for measuring non-linear refraction and absorption.
    \item To model intensity-dependent refractive indexes using python-based numerical simulations.
    \item To verify simulation runs against experimental values of typical optical glasses.
\end{itemize}

\section{Scope and Organization of the Report}
This project report is organized as follows: Chapter \ref{ch:methodology} details the background theory and experimental methodology. Chapter \ref{ch:results} outlines the results and plots. Chapter \ref{ch:discussion} provides interpretations and comparison, while Chapter \ref{ch:conclusion} concludes the project report and outlines potential future improvements.
